% ============================================================
% PAGE 3 — DEEP LEARNING PARADIGM SHIFT & TWO-STAGE DETECTORS PART 1
% IEEE two-column content (Overleaf-ready)
% Paste below previous pages and before \bibliographystyle / \bibliography
% ============================================================

% ============================================================
% SECTION 3: THE DEEP LEARNING PARADIGM SHIFT
% ============================================================

\section{The Deep Learning Paradigm Shift}

\subsection{AlexNet and the Redefinition of Feature Learning}

The year 2012 marked a decisive turning point in computer vision.
Prior to this, object detection and recognition pipelines relied
on features designed by hand: Histogram of Oriented Gradients
(HOG) encoded local intensity gradients; Scale-Invariant Feature
Transform (SIFT) captured distinctive key-point descriptors;
and bag-of-words models aggregated these descriptors into global
image representations. Each of these operators embodied
domain-specific human intuition about what visual information
matters, and improving a pipeline meant refining that intuition
through careful, manual iteration.

AlexNet~\cite{krizhevsky2012alexnet}, a deep Convolutional Neural
Network (CNN) submitted by Krizhevsky, Sutskever, and Hinton to
the ImageNet Large Scale Visual Recognition Challenge (ILSVRC)
2012, achieved a top-5 error rate of 15.3\%---more than ten
percentage points lower than the next best entry. The architecture
stacked five convolutional layers followed by three fully connected
layers, trained on approximately 1.2 million labeled images using
stochastic gradient descent with momentum. Several design choices
proved critical: the Rectified Linear Unit (ReLU) activation
function accelerated convergence compared to saturating
activations; dropout regularization reduced over-fitting in the
fully connected layers; and GPU-based training made the otherwise
prohibitive parameter count tractable.

The broader implication was architectural rather than incremental.
AlexNet demonstrated that a sufficiently deep network, trained
with sufficient data and compute, \emph{learns} which visual
abstractions are discriminative---without the engineer specifying
them explicitly. Early convolutional layers were observed to
activate on low-level primitives (oriented edges, color blobs),
intermediate layers on textures and object parts, and deeper
layers on semantically coherent patterns. This hierarchy emerged
from the training signal alone.

\subsection{From Hand-Crafted to Learned Representations}

The transition from engineered to learned features transformed
the object detection problem in two important ways.

\textbf{Representation quality.} Features extracted from deep
CNN layers are substantially richer in discriminative content
than hand-crafted alternatives. A deep feature vector encodes
not just the presence of edges but also the semantic context
in which they appear---an edge at a wheel arch activates
differently from an edge on a building corner. This semantic
grounding is precisely what hand-crafted features, operating
locally and without end-to-end optimization, cannot capture.

\textbf{End-to-end optimization.} Because CNN weights are learned
jointly with the classification objective, the representation
is optimized for the task rather than for some proxy criterion.
When the detection loss propagates gradients back through the
network, every layer adjusts its parameters to minimize
classification and localization error simultaneously---a form
of task-aware feature learning unavailable in fixed-feature
pipelines.

The success of AlexNet on image classification raised an immediate
question for the object detection community: can the same deep
features be applied to detection, where multiple objects appear
at varying positions and scales within a single image? This
question prompted a wave of research that sought to adapt CNN-based
classification to the structured, multi-object detection setting,
giving rise to the first generation of deep learning-based
detectors.

% ============================================================
% SECTION 4: TWO-STAGE DETECTORS — PART 1
% ============================================================

\section{Two-Stage Detectors: Part~1}

\subsection{The Region Proposal Principle}

Applying a CNN classifier naively to detection---by sliding a
fixed window across every position and scale in an image---is
computationally intractable. For a $600 \times 800$ image and
even a modest set of aspect ratios and scales, thousands of
candidate windows must be evaluated, each requiring a full
forward pass through a deep network. This naive brute-force
approach does not scale.

The two-stage detection paradigm resolves this by decomposing
the problem into two sequential steps:

\begin{enumerate}
    \item \textbf{Region proposal:} Generate a modest set of
    candidate bounding boxes (regions of interest) that are
    likely to contain objects, without yet classifying them.
    \item \textbf{Region classification:} Apply a more powerful
    classifier to each candidate region independently, producing
    final class scores and refined box coordinates.
\end{enumerate}

This decomposition trades off the exhaustive coverage of sliding
windows for a principled, data-driven reduction in the number of
candidates that need detailed analysis. If the proposal step
generates a few hundred to a few thousand candidates rather than
millions, the per-region classification cost becomes manageable---
while overall recall can remain high provided the proposal
mechanism is sufficiently sensitive.

\subsection{R-CNN: Regions with CNN Features}

Girshick et al.~\cite{girshick2014rcnn} introduced R-CNN
(Regions with Convolutional Neural Network Features) as the first
detector to combine selective region proposals with deep CNN
classification. The pipeline consists of three stages.

\textbf{Selective Search proposals.}
Rather than an exhaustive grid, R-CNN uses Selective
Search~\cite{uijlings2013selective} to generate approximately
2,000 candidate regions per image. Selective Search is a
bottom-up grouping algorithm that merges superpixels based on
color, texture, size, and fill similarity, producing a
hierarchical set of regions at multiple scales. Because it
operates on the image directly and exploits pixel-level
grouping cues, it tends to produce regions that respect object
boundaries---yielding high recall with far fewer candidates than
a dense sliding window.

\textbf{Per-region CNN feature extraction.}
Each candidate region is warped to a fixed $227 \times 227$ pixel
crop and passed independently through a pre-trained CNN (AlexNet
in the original paper, later VGG-16 in follow-up experiments)
to extract a fixed-length feature vector. Pre-training on
ImageNet and then fine-tuning on PASCAL VOC ground-truth boxes
allows the CNN to adapt its representations toward the detection
objective.

\textbf{SVM classification and bounding-box regression.}
The extracted CNN feature vector is fed into a set of
class-specific linear Support Vector Machines (SVMs) to score
each candidate region. A separate linear bounding-box regressor
refines the proposed coordinates to better align with ground
truth. Because the CNN feature extractor and the SVM classifier
are trained separately, the overall pipeline is a multi-stage
system with several independent training steps.

R-CNN substantially outperformed HOG-SVM baselines on PASCAL
VOC 2012, demonstrating that learned CNN features transfer
effectively to detection. However, its architecture imposes a
significant computational cost: because each of the $\sim$2,000
proposals requires an independent CNN forward pass, processing
a single image takes approximately 49 seconds on a GPU at test
time. The three-stage training procedure---CNN pre-training,
fine-tuning, SVM fitting, and bounding-box regressor fitting---
also complicates practical use.

\subsection{Fast R-CNN: Shared Convolutional Features}

Girshick~\cite{girshick2015fastrcnn} addressed the efficiency
limitations of R-CNN in Fast R-CNN by fundamentally reordering
the computation: instead of warping and encoding each region
proposal independently, Fast R-CNN passes the \emph{entire
image} through the convolutional backbone \emph{once},
producing a shared convolutional feature map, and then projects
each proposal onto that feature map.

\textbf{Region of Interest (RoI) pooling.}
Each candidate region, originally defined in image coordinates,
is mapped to a corresponding sub-region of the convolutional
feature map. A RoI pooling layer divides this sub-region into
a fixed $H \times W$ grid and applies max-pooling within each
cell, producing a fixed-size feature vector regardless of the
original region's aspect ratio or size. Because all proposals
share the same underlying feature map, the expensive
convolutional computation is performed only once per image
rather than once per region.

\textbf{Multi-task loss and end-to-end training.}
Fast R-CNN replaces the separate SVM and regressor training steps
of R-CNN with a single, jointly-optimized objective. A softmax
classifier produces class probabilities, and a class-specific
bounding-box regression head predicts coordinate offsets. Both
are trained simultaneously under a multi-task loss:

\begin{equation}
    \mathcal{L}(p, u, t^{u}, v) =
    \mathcal{L}_{\text{cls}}(p, u)
    + \lambda \, [u \geq 1] \, \mathcal{L}_{\text{loc}}(t^{u}, v),
    \label{eq:fast_rcnn_loss}
\end{equation}

\noindent
where $p = (p_0, \ldots, p_K)$ is the discrete probability
distribution over $K+1$ classes, $u$ is the ground-truth class,
$t^u$ is the predicted offsets for class $u$, $v$ is the
ground-truth box, and $\lambda$ balances the two terms.
The indicator $[u \geq 1]$ suppresses the localization loss for
background regions. $\mathcal{L}_{\text{loc}}$ is the smooth-$L_1$
loss, which is less sensitive to outlier boxes than the squared
error.

\textbf{Speed improvement over R-CNN.}
By sharing the convolutional backbone across all proposals, Fast
R-CNN reduces test-time inference from $\sim$49 seconds per image
(R-CNN with VGG-16) to $\sim$0.3 seconds---an approximately
160$\times$ improvement---while simultaneously improving mAP on
PASCAL VOC 2012 from 62.4\% to 68.4\%. End-to-end training
further contributes to this accuracy gain by allowing the
convolutional layers to be optimized for the detection objective
rather than being held fixed. The primary remaining bottleneck
in Fast R-CNN is the Selective Search proposal step, which runs
on the CPU and accounts for the majority of remaining inference
time---a limitation that motivated the subsequent development of
learned, network-integrated proposal mechanisms.

% ============================================================
% BIBLIOGRAPHY ADDITIONS (add these entries to your .bib file)
% ============================================================
%
% @inproceedings{krizhevsky2012alexnet,
%   author    = {Krizhevsky, Alex and Sutskever, Ilya and Hinton, Geoffrey E.},
%   title     = {{ImageNet} Classification with Deep Convolutional Neural Networks},
%   booktitle = {Advances in Neural Information Processing Systems (NeurIPS)},
%   volume    = {25},
%   year      = {2012}
% }
%
% @inproceedings{girshick2014rcnn,
%   author    = {Girshick, Ross and Donahue, Jeff and Darrell, Trevor and Malik, Jitendra},
%   title     = {Rich Feature Hierarchies for Accurate Object Detection and Semantic Segmentation},
%   booktitle = {Proceedings of the IEEE Conference on Computer Vision and Pattern Recognition (CVPR)},
%   year      = {2014},
%   pages     = {580--587}
% }
%
% @article{uijlings2013selective,
%   author  = {Uijlings, Jasper R. R. and van de Sande, Koen E. A. and Gevers, Theo and Smeulders, Arnold W. M.},
%   title   = {Selective Search for Object Recognition},
%   journal = {International Journal of Computer Vision (IJCV)},
%   volume  = {104},
%   number  = {2},
%   pages   = {154--171},
%   year    = {2013}
% }
%
% @inproceedings{girshick2015fastrcnn,
%   author    = {Girshick, Ross},
%   title     = {Fast {R-CNN}},
%   booktitle = {Proceedings of the IEEE International Conference on Computer Vision (ICCV)},
%   year      = {2015},
%   pages     = {1440--1448}
% }
